\section{Abelian varieties of Weil type}\label{sec:weiltype}

This section is self-contained. We fix the notion of Weil type, construct the
Weil line, prove that it consists of Hodge classes, and record the two facts
about it that the obstruction theorems use: its dimension, and its position
relative to the subring generated by the polarisation. Standard references for
the background are \cite{BL04} for abelian varieties and \cite{Voi02} for Hodge
theory; nothing below is quoted without proof except the Riemann relations and
the Lefschetz theorem on $(1,1)$ classes.

\subsection{Weil type}

Let $K=\QQ(\sqrt{-d})$ be an imaginary quadratic field, $d$ a squarefree
positive integer, and let $\bbar{\phantom{x}}$ denote the nontrivial element
of $\Gal(K/\QQ)$.

\begin{definition}[Weil type]\label{def:weiltype}
Let $A$ be a complex abelian variety of dimension $2n$, let
$\iota\colon K\hookrightarrow\End^{0}(A)$ be an embedding, and let $\eta$ be a
polarisation class on $A$. The triple $(A,\iota,\eta)$ is of
\emph{$(K,-1,n)$-Weil type} if%
\footnote{The middle entry records the sign in the compatibility between the polarisation and the $K$-action, and is $-1$ here because the Rosati involution restricts to conjugation rather than to the identity. Some authors suppress it and write simply $K$-Weil type of dimension $2n$; nothing below depends on the convention, and by \Cref{lem:rosatiauto} the sign is in any case forced.}
\begin{enumerate}[label=\textup{(\roman*)},nosep]
\item $\dim A=2n$;
\item the endomorphism $\iota(\sqrt{-d})$ acts on the tangent space
  $H^{0}(A,\Omega^{1}_{A})^{\vee}\cong H^{1,0}(A)^{\vee}$ with the two
  eigenvalues $\pm\sqrt{-d}$ each occurring with multiplicity $n$;
\item the Rosati involution attached to $\eta$ restricts on
  $\iota(K)\subset\End^{0}(A)$ to complex conjugation, that is
  $\iota(\tau)^{\dagger}=\iota(\bbar{\tau})$ for all $\tau\in K$.
\end{enumerate}
\end{definition}

Condition (ii) is the balanced signature condition; it is what distinguishes
Weil type from the generic situation, and it is what makes the Weil line
consist of Hodge classes. Condition (iii) says that the polarisation is
compatible with the $K$-action.

\begin{lemma}[$H^{1}$ is a $K$-space]\label{lem:Kspace}
Let $(A,\iota,\eta)$ be of $(K,-1,n)$-Weil type and put $V=H^{1}(A,\QQ)$. Then
$\iota$ makes $V$ a $K$-vector space of dimension $2n$, and
$\dim_{\QQ}V=4n$.
\end{lemma}

\begin{proof}
The map $\iota$ embeds $K$ into $\End^{0}(A)$, and $\End^{0}(A)$ acts
$\QQ$-linearly on $V$ by pullback. Since $K$ is a field, $V$ becomes a
$K$-vector space, and $\dim_{\QQ}V=2\dim A=4n$, so $\dim_{K}V=4n/[K:\QQ]=2n$.
\end{proof}

\begin{notation}\label{not:eigen}
Write $V_{\CC}=V\otimes_{\QQ}\CC$. Since $\sqrt{-d}$ acts on $V$ with
$(\sqrt{-d})^{2}=-d$, the operator $\iota(\sqrt{-d})\otimes1$ on $V_{\CC}$ has
the two eigenvalues $\pm\sqrt{-d}$, and
\begin{equation}\label{eq:Vsplit}
  V_{\CC}=V_{+}\oplus V_{-},
  \qquad
  V_{\pm}=\ker\bigl(\iota(\sqrt{-d})\mp\sqrt{-d}\bigr),
\end{equation}
with $\dim_{\CC}V_{\pm}=2n$ each, because $V$ is free of rank $2n$ over $K$ and
$K\otimes_{\QQ}\CC\cong\CC\times\CC$. Complex conjugation on $V_{\CC}$, coming
from the real structure $V\otimes\RR$, interchanges $V_{+}$ and $V_{-}$. For
general $\tau=a+b\sqrt{-d}\in K$ the operator $\iota(\tau)$ acts on $V_{+}$ by
the scalar $\tau$ and on $V_{-}$ by $\bbar{\tau}$.
\end{notation}

\begin{remark}\label{rem:tautbar}
The two scalars in \Cref{not:eigen} are different functions of $\tau$: the operator $\iota(\tau)$ is
$\QQ$-linear on $V$ but not a scalar, and only after extending to $\CC$ and
splitting does it become scalar on each piece. This is the whole source of the
character phenomenon in \Cref{sec:characters}. A construction defined over
$\QQ$ sees only the $\QQ$-linear operator, and hence only those of its
eigenvalues that are $\QQ$-rational functions of $\tau$; $\Nm(\tau)=\tau\bbar{\tau}$
is one such, and $\tau^{2n}$ is not.
\end{remark}

\subsection{The Hodge structure}

The Hodge decomposition of $V_{\CC}$ is $V_{\CC}=H^{1,0}\oplus H^{0,1}$, and
the $K$-action preserves it because $\iota(\tau)$ is induced by an
endomorphism of $A$ and hence is a morphism of Hodge structures. Therefore
each of $V_{+}$ and $V_{-}$ splits,
\begin{equation}\label{eq:bisplit}
  V_{\pm}=V_{\pm}^{1,0}\oplus V_{\pm}^{0,1}.
\end{equation}

\begin{lemma}[Balanced signature]\label{lem:balanced}
For $(A,\iota,\eta)$ of $(K,-1,n)$-Weil type,
\[
  \dim_{\CC}V_{+}^{1,0}=\dim_{\CC}V_{-}^{1,0}=n,
  \qquad
  \dim_{\CC}V_{+}^{0,1}=\dim_{\CC}V_{-}^{0,1}=n .
\]
\end{lemma}

\begin{proof}
Condition (ii) of \Cref{def:weiltype} says that $\iota(\sqrt{-d})$ acts on the
tangent space with each eigenvalue of multiplicity $n$. The tangent space is
dual to $H^{1,0}(A)=V^{1,0}$, so $\iota(\sqrt{-d})$ acts on $V^{1,0}$ with each
eigenvalue of multiplicity $n$; that is, $\dim V^{1,0}_{+}=\dim V^{1,0}_{-}=n$.
Since $\dim_{\CC}V_{\pm}=2n$ by \Cref{not:eigen} and
$V_{\pm}=V^{1,0}_{\pm}\oplus V^{0,1}_{\pm}$, the remaining two dimensions are
$n$ each.
\end{proof}

\subsection{The Weil line}

\begin{definition}[Weil line]\label{def:weilline}
Let $(A,\iota,\eta)$ be of $(K,-1,n)$-Weil type and $V=H^{1}(A,\QQ)$. The
\emph{Weil line} is
\[
  \HW(A)\;=\;\Bigl(\textstyle\bigwedge^{2n}_{\CC}V_{+}\oplus
  \bigwedge^{2n}_{\CC}V_{-}\Bigr)\cap H^{2n}(A,\QQ),
\]
the intersection being taken inside
$H^{2n}(A,\CC)=\bigwedge^{2n}_{\CC}V_{\CC}$.
\end{definition}

\begin{proposition}[Structure of the Weil line]\label{prop:weilline}
Let $(A,\iota,\eta)$ be of $(K,-1,n)$-Weil type. Then
\begin{enumerate}[label=\textup{(\roman*)},nosep]
\item $\HW(A)$ is a $\QQ$-subspace of $H^{2n}(A,\QQ)$ of dimension $2$, stable
  under $\iota(K)$, and free of rank one as a $K$-module;
\item $\HW(A)\subset H^{n,n}(A)$, so every element of $\HW(A)$ is a Hodge class;
\item $\HW(A)$ is canonically isomorphic to $\bigwedge^{2n}_{K}V$ as a
  $K$-module.
\end{enumerate}
\end{proposition}

\begin{proof}
(i) Write $W_{\CC}=\bigwedge^{2n}V_{+}\oplus\bigwedge^{2n}V_{-}$. Each summand
is one-dimensional over $\CC$ because $\dim_{\CC}V_{\pm}=2n$, so
$\dim_{\CC}W_{\CC}=2$. Complex conjugation on $H^{2n}(A,\CC)$ interchanges
$V_{+}$ and $V_{-}$ and hence interchanges the two summands, so $W_{\CC}$ is
defined over $\RR$; it is in fact defined over $\QQ$, because the splitting
\eqref{eq:Vsplit} is the eigenspace decomposition of an operator
$\iota(\sqrt{-d})$ defined over $\QQ$, and the pair of extreme exterior powers
is permuted by $\Gal(\CC/\QQ)$ acting through $\Gal(K/\QQ)$. Therefore
$\HW(A)=W_{\CC}\cap H^{2n}(A,\QQ)$ has $\dim_{\QQ}\HW(A)=2$ and
$\HW(A)\otimes\CC=W_{\CC}$. Stability under $\iota(K)$ is clear since
$\iota(\tau)$ preserves each of $V_{\pm}$, hence each summand. A
two-dimensional $\QQ$-space with a faithful $K$-action is free of rank one
over $K$.

(ii) The Hodge decomposition of $\bigwedge^{2n}V_{+}$ is computed from
\eqref{eq:bisplit}: the single line $\bigwedge^{2n}V_{+}$ equals
$\bigwedge^{n}V^{1,0}_{+}\otimes\bigwedge^{n}V^{0,1}_{+}$, which by
\Cref{lem:balanced} is of Hodge bidegree $(n,n)$. The same computation applies
to $\bigwedge^{2n}V_{-}$. Hence $W_{\CC}\subset H^{n,n}(A)$ and so
$\HW(A)\subset H^{2n}(A,\QQ)\cap H^{n,n}(A)=\Hdg^{n}(A)$.

(iii) The natural map $\bigwedge^{2n}_{K}V\to\bigwedge^{2n}_{\QQ}V$ obtained by
antisymmetrising over a $K$-basis identifies $\bigwedge^{2n}_{K}V$ with a
$K$-line in $H^{2n}(A,\QQ)$, and after $\otimes\CC$ the $K$-structure splits
this line into the $\tau$-eigenline and the $\bbar{\tau}$-eigenline, which are
$\bigwedge^{2n}V_{+}$ and $\bigwedge^{2n}V_{-}$ respectively. Both sides are
two-dimensional over $\QQ$ and the map is injective, hence an isomorphism.
\end{proof}

\begin{remark}\label{rem:onedim}
Some authors call a generator of $\HW(A)$ \emph{the} Weil class and speak of a
one-dimensional space; the discrepancy is the usual one between a $K$-line and
its underlying $\QQ$-plane. Nothing below depends on the convention, because
both obstruction theorems are statements about the whole of $\HW(A)$.
\end{remark}

\subsection{Weil classes are not polynomials in the polarisation}

The point of the Weil line is that it is not accounted for by divisors. We
record the statement in the form we use and prove it from the character
computation of the next section, which is the cleanest route.

\begin{definition}\label{def:divisorsubring}
The \emph{divisor subring} $D^{*}(A)\subset H^{*}(A,\QQ)$ is the
$\QQ$-subalgebra generated by $\NS(A)\otimes\QQ\subset H^{2}(A,\QQ)$.
\end{definition}

Every element of $D^{*}(A)$ is algebraic, by the Lefschetz theorem on $(1,1)$
classes together with the fact that the algebraic classes form a subring. The
Hodge conjecture for $A$ in degree $n$ is therefore equivalent to the
algebraicity of a complement of $D^{2n}(A)$ in $\Hdg^{n}(A)$, and by the
structure theory of Hodge classes on abelian varieties of Weil type the Weil
line is that complement.

\begin{proposition}\label{prop:notpoly}
Let $(A,\iota,\eta)$ be of $(K,-1,n)$-Weil type with
$\End^{0}(A)=\iota(K)$, and suppose $n\ge1$. Then
\[
  \HW(A)\cap D^{2n}(A)=0 .
\]
\end{proposition}

\begin{proof}
By \Cref{prop:weilline}(i) and \Cref{thm:character}, proved independently in
\Cref{sec:characters}, the space $\HW(A)\otimes\CC$ is the sum of the
$\Kx$-eigenspaces for the characters $\tau\mapsto\tau^{2n}$ and
$\tau\mapsto\bbar{\tau}^{2n}$. By \Cref{lem:NSchar} below, every class in
$\NS(A)\otimes\QQ$ is a $\Kx$-eigenclass for the norm character $\Nm$, so
every homogeneous element of $D^{2n}(A)$ is an eigenclass for $\Nm^{n}$.%
\footnote{This is where the whole argument sits, so we put it in words. A divisor class is a rank one object, so the field acts on it through a quantity a rank one object can see, and over $\QQ$ that quantity is the norm. The Weil line is not of rank one over $\QQ$; it is of rank one over $K$, and a $K$-line of exterior degree $2n$ carries the $2n$-th power of $\tau$ itself.} By
\Cref{lem:charsdiffer} the characters $\tau\mapsto\tau^{2n}$ and
$\tau\mapsto\Nm(\tau)^{n}$ are distinct, and eigenspaces for distinct
characters intersect in zero.
\end{proof}

We prove \Cref{lem:NSchar} and \Cref{lem:charsdiffer} in the next section; the
forward reference is harmless because nothing in \Cref{sec:characters} uses
\Cref{prop:notpoly}.
